\newcommand{\TriptychBriefSynthesis}{%
\section*{The Propers: Themes and Movement}
\label{triptych:brief-synthesis:start}

\begin{massbox}{Presumed standing yields to received abiding}
God helps the threatened and reforms desire; Paul turns sacramental privilege into a warning against self-reliance; Christ visits Jerusalem through tears and prophetic correction; the Eucharistic memorial gives the abiding, purification, and one-body unity that human standing cannot secure.
\end{massbox}

\begin{enumerate}[leftmargin=1.5em,itemsep=0.1em]
\item \textbf{Receive help:} mercy reforms desire.
\item \textbf{Relinquish presumption:} figure becomes warning.
\item \textbf{Receive visitation:} tears issue in reform and teaching.
\item \textbf{Abide as one:} memorial bears fruit in a body.
\end{enumerate}

\noindent Read in these four stages, the appointed sequence moves from received help to abiding as one, and the witnesses cited within each stage ground it.

\subsection*{Help reaches the petitioner by reforming desire\enspace\properrefs{Int., Coll., All.}}

Psalm 53's threatened speaker discovers God already helping and sustaining, yet continues to ask rescue. The Collect carries that dependence inward: before desired things are granted, God must teach the petitioners to seek what pleases him. The Alleluia again entrusts deliverance to God. Mercy neither ratifies every impulse nor removes responsible action; it reforms the one who asks.

The three appointed texts keep distinct moments of dependence together. The Introit's ``behold'' recognizes a present helper while danger remains; the Collect submits the petitioner's own desiring to formation; the Alleluia returns to an urgent imperative for rescue. Confidence therefore neither pretends that the threat has vanished nor waits passively. Johner's chant commentary hears the Introit move from firm confidence through unrest and back to confidence, while the Alleluia's plea remains enclosed by praise.

\subsection*{Privilege becomes warning to those who think they stand\enspace\properrefs{Ep., Grad.}}

Paul first recalls sea, cloud, food, and drink in sacramental language, then makes Israel's wilderness failures correction for the Corinthian baptized. Craving, idolatry, sexual sin, testing, murmuring, and the claim to stand expose sacramental presumption. Psalm 8 answers with creaturely praise beneath God's world-filling name: dignity is received, not self-secured. God's fidelity supplies endurance within temptation.

Ambrosiaster sharpens the social force of ``thinks he stands.'' He applies it to Corinthians whose asserted knowledge and liberty around idol food scandalized weaker Christians: what they counted as advance had made them worse. Standing is thus tested by what freedom does to another member. His exposition then follows Paul from the promised capacity to endure into flight from idolatry and the one bread and body.

\begin{studybox}{The canonical hinge beyond the printed lesson}
The missal stops at 1~Corinthians 10:13, but Paul's argument does not. ``Therefore'' leads to flight from idolatry, then to the Eucharistic cup and bread and the conclusion that the many are one body because they share one bread (10:14--17). The Sunday later sets the self-estimate \latin{qui se existimat stare} beside Communion's received relation \latin{in me manet, et ego in eo}. That contrast and the movement to unity are consonant with Paul's argument; the arrangement that sets them together is Roman.
\end{studybox}

\newpage

\subsection*{Visitation arrives as tears, correction, and teaching\enspace\properrefs{Gosp., Off.}}

Luke's royal approach proclaims peace; Jesus then weeps over peace unrecognized, foretells the city's enclosure, enters the Temple, and teaches there daily. Isaiah's ``house of prayer'' opens worship to all peoples; Jeremiah's ``den of robbers'' condemns trust in sanctuary without justice toward neighbor, alien, orphan, and widow. The action is Jewish prophetic correction ordered to prayer and teaching. In the appointed Psalm 18 span, right precepts gladden, the command enlightens, true judgments are desired, and the servant is warned and keeps them. Together the Gospel and Offertory direct judgment toward prayer, teaching, delight in truth, warning, and obedient keeping.

Origen's \emph{Homily 38 on Luke} sees Christ's tears enact the beatitude of mourning and then directs the warning toward Christians who have received Gospel teaching and the sacraments. The Pontifical Biblical Commission likewise places the tears beside Luke 23:28--31, where Jesus attends to Jerusalem's women and children rather than his own suffering. Ambrose turns the Temple action toward Christian ministry: holiness is not merchandise, Scripture is a divine coin to be discerned, and grace received through the Spirit is not sold. These applications call Christian ministry to examination while Christ's tears keep Jerusalem's suffering in view.

\subsection*{The memorial gives abiding and tests it by unity\enspace\properrefs{Sec., Comm., Postcomm.}}

The Secret says that whenever the sacrifice's memorial is celebrated, redemption's work is enacted. Thomas directly quotes this Secret while distinguishing the once-for-all Passion from daily participation in its fruit; Vatican II twice identifies the prayer as its source. \emph{Lumen gentium} immediately joins that action to Eucharistic unity in 1~Corinthians 10:17. John names the gift as mutual abiding; the Postcommunion asks its fruit as purification and unity.

Origen ends the same Gospel homily by summoning hearers to the Savior's banquet, the bread of life, and Christ's flesh and blood. Bede follows his Luke 19 exposition with a Paschal synthesis that moves from Christ's reception and Temple teaching through the gift of his Body and Blood to the Passion. Origen and Bede thus cross from visitation to Eucharistic reception within their expositions of Luke. Dom Lucien Fromage's direct commentary on this Sunday, in the continuation of Guéranger's \emph{The Liturgical Year}, then calls Communion divine union and identifies the Postcommunion's paired fruits as personal sanctification and the unity of the ecclesial body. The Trinity Preface directs the whole reception toward praise of the triune God.

\begin{threestable}{Movement}{Decisive gift}{Precise sense}
\textbf{Help received} & God sustains, reforms desire, and hears the endangered petitioner. & Mercy neither endorses every desire nor abolishes responsible action.\\
\textbf{Standing chastened} & God's fidelity makes endurance possible and turns the baptized from idolatry. & Verse 13 concerns temptation, not every trauma or isolated human capacity.\\
\textbf{Visitation received} & Christ weeps, judges false worship, and remains to teach. & Tears, reform, and continued teaching remain one narrated movement.\\
\textbf{Abiding embodied} & The memorial enacts redemption; received Communion is asked to bear purification and unity amid triune praise. & The Passion is not repeated, and sacramental fruit is not private possession.\\
\end{threestable}

\label{triptych:brief-synthesis:end}
\clearpage
}
